\section{Spiritual Attitudes}

\begin{quotex}
One must live as one thinks, under pain of sooner or later thinking as one lives.
\flright{\textsc{Paul Bourget}}
\end{quotex}


In \emph{The Individual and the becoming of the World}, Evola describes certain spiritual ``attitudes'', a concept that will be developed and expanded in his later works as ``spiritual types'' and then ``races of the spirit''. The editors of \emph{Sintesi di dottrina della razza} define it thus: ``The `race of the spirit' prevalent in a person, a people, a community is given by the characteristic orientation that is assumed toward the sacred and the divine, life and death, destiny, the world.''

We moderns, under the influence of universalism and egalitarianism, tend to underestimate the significances of spiritual types: Don't we all laugh and cry? Don't we bleed when pricked? These, and other such rhetorical questions, seek the unity of humanity in the lowest and most common features. In actuality, men of different spiritual types, which are far from arbitrary, represent quite different states of being. This is obvious when men of differing spiritual types attempt a conversation -- their disparate fundamental orientations preclude a common understanding of anything above the most superficial topics.

This needs to be clarified. The lower man cannot understand the higher. However, as Weininger points out, the higher man can understand the lower since he encompasses more. Therefore, the man of spirit can see that the materialistic may be correct \emph{as far as it goes} ... he simply does not go far enough. However, the materialist can only regard the spiritual man as mistaken, if not deluded.

It is important to keep in mind that the spirit, as principle, is primary. Its characteristics are then manifested in the soul and the body. In the chapter ``The Illusion of Ordinary Life'' in \emph{The Reign of Quantity \& the Signs of the Time}, Guenon writes this:

\begin{quotex}
The materialistic attitude, whether it be a question of explicit and formal materialism or of a simple ``practical'' materialism, necessarily imposed on the whole ``psycho-physiological'' constitution of the human being a real and very important modification. This is easily understood, and in fact it is only necessary to look around in order to conclude that modern man has become quite impermeable to any influences other than such as impinge on his senses; not only have his faculties of comprehension becomes more and more limited, but also the field of his perception has become correspondingly restricted. The result is a sort of reinforcement of the profane point of view, for this point of view was first born of a defect of comprehension, and thus of a limitation, and this limitation, as it becomes accentuated and extends to all domains, itself seems to justify the point of view, at least in the eyes of those who are affected by it. Indeed, what reason can they have thereafter for admitting the existence of something that they can neither perceive not conceive, that is to say of everything that could show them the insufficiency and the falsity of the profane point of view itself?
\end{quotex}


In \emph{Meditations on the Peaks}, Evola is explicit:

\begin{quotex}
All too often people forget that spirituality is essentially a way of life and that its measure does not consist of notions, theories, and ideals that have been stored in one's head. Spirituality is actually what has been successfully actualized and translated into a sense of superiority which is experienced inside by the soul, and a noble demeanor, which is expressed in the body.
\end{quotex}


So in the second part of ``The Individual'', Evola does not at all concern himself with notions, theories, and ideals. Rather, he is concerned about which spiritual attitudes lead to that inner sense of superiority and which fall short.


\flrightit{Posted on 2008-09-15 by Cologero}